\section{Methods}
\label{sec:methods}

\subsection{GITT processing}
We process galvanostatic intermittent titration technique (GITT) measurements
to compute defensible, step-level metrics such as rest voltage change
($\Delta E_s$), pulse voltage change ($\Delta E_\tau$), pulse duration ($\tau$),
and early-time $V$--$\sqrt{t}$ slope. Following classic treatments of GITT
\citep{WeppnerHuggins1977GITT}, an \emph{apparent} diffusion coefficient may be
computed only if an explicit diffusion length scale $L$ is assumed and stated.
Because full-cell voltage contains contributions from both electrodes and
overpotentials, this work does not claim electrode-specific solid diffusivities
from full-cell GITT alone.

\subsection{Physics model: \pybamm}
We use \pybamm \citep{Sulzer2021PyBaMM} to define an electrochemical model
family for LFP/graphite cells. Within this repository, thermal effects are
disabled (isothermal at 25\,$^\circ$C) to match the dataset conditions.

The role of \pybamm is to provide a physics prior and synthetic trajectory
generator. To be scientifically defensible, electrochemical and degradation
parameters must be calibrated only against tests that actually constrain them,
and validated on held-out protocols within the same operating envelope.

\subsection{Neural ODE surrogate for \SOH dynamics}
We model \SOH evolution across cycle number $n$ via a continuous-depth ODE:
\begin{equation}
  \frac{d\,\SOH}{dn} = f_\theta(\SOH, n, \mathbf{x}),
\end{equation}
where $\mathbf{x}$ is a vector of health features (e.g., DCIR and IC features).
We integrate this system using an adaptive Runge--Kutta method as implemented in
\texttt{torchdiffeq} \citep{Chen2018NeuralODE}.

\paragraph{Monotonicity constraint.}
To encode the physical prior that capacity-related \SOH should not increase
under normal aging, we enforce $d\SOH/dn \le 0$ by parameterising the ODE output
as $-\mathrm{softplus}(\cdot)$ in the implementation.

\paragraph{Relationship to \pinn terminology.}
The codebase refers to this as a \pinn-style approach in the broad sense of
embedding physics constraints into training. However, because the governing
system is an ODE in cycle count rather than a PDE in space--time, the surrogate
is best described as a physics-guided Neural ODE with monotonic constraints,
and any additional ``physics residual'' terms must be justified against
validated mechanistic models.

